% Part: turing-machines 
% Chapter: undecidability
% Section: introduction

\documentclass[../../../include/open-logic-section]{subfiles}

\begin{document}

\olfileid[hi]{tur}{und}{int} 
\olsection{परिचय}

यह सहज लग सकता है कि प्रत्येक फलन---यहाँ तक कि प्रत्येक अंकगणितीय
फलन भी---संगणनीय नहीं हो सकता। ऐसे फलन अत्यधिक हैं और उनका व्यवहार
बहुत जटिल है। उदाहरणतः रेडियोधर्मी कणों के क्षय, अथवा अन्य अराजक या
यादृच्छिक व्यवहार से परिभाषित फलन। मान लें कि हम किसी निश्चित समय से
एक-सेकंड के अंतराल गिनना आरंभ करते हैं और फलन $f(n)$ को उस आरंभिक क्षण
के बाद $n$-वें एक-सेकंड अंतराल में ब्रह्मांड में क्षय होने वाले कणों की
संख्या के रूप में परिभाषित करते हैं। यह ऐसे फलन का प्रत्याशी लगता है
जिसकी संगणना कर पाने की आशा हम कभी नहीं कर सकते।

किंतु ऐसे फलनों की संगणना की विधि की कल्पना न कर पाना एक बात है और
वास्तव में यह सिद्ध करना बिल्कुल दूसरी बात कि वे असंगणनीय हैं। वस्तुतः
जो फलन निराशाजनक रूप से जटिल लगते हैं, वे भी अमूर्त अर्थ में संगणनीय हो
सकते हैं। उदाहरण के लिए मान लें कि ब्रह्मांड समय की दृष्टि से परिमित
है---कुछ ब्रह्मांड-विज्ञान संबंधी सिद्धांतों के पूर्वानुमान के अनुसार
बहुत दूर भविष्य में किसी दिन ब्रह्मांड सिकुड़कर एक बिंदु बन जाएगा। तब
उस आरंभिक क्षण से केवल परिमित (यद्यपि अविश्वसनीय रूप से विशाल) संख्या
में ऐसे सेकंड होंगे जिनके लिए $f(n)$ परिभाषित है। और केवल परिमित संख्या
में आगतों पर परिभाषित प्रत्येक फलन संगणनीय है: हम उसके निर्गतों को एक
बड़ी सारणी में लिख सकते हैं, अथवा उसे ट्यूरिंग मशीन के एक बहुत बड़े
अवस्था-संक्रमण आरेख में कूटबद्ध कर सकते हैं।

हम प्रायः फलनों के उन विशेष मामलों में रुचि रखते हैं जिनके मान हाँ/नहीं
प्रश्नों के उत्तर देते हैं। उदाहरणतः प्रश्न ``क्या $n$ एक अभाज्य संख्या
है?'' निम्न फलन से संबद्ध है:
\[
\fn{isprime}(n) = \begin{cases}
  1 & \text{यदि $n$ अभाज्य है}\\
  0 & \text{अन्यथा।}
  \end{cases}
\]
हम कहते हैं कि किसी हाँ/नहीं प्रश्न का \emph{प्रभावी रूप से निर्णय} किया
जा सकता है, यदि उससे संबद्ध $1/0$-मान वाला फलन प्रभावी रूप से संगणनीय
हो।

गणितीय रूप से यह सिद्ध करने के लिए कि कुछ फलन प्रभावी रूप से संगणित
नहीं किए जा सकते अथवा कुछ समस्याओं का प्रभावी रूप से निर्णय नहीं किया
जा सकता, संगणना का कोई विशिष्ट प्रतिरूप निश्चित करना और यह दिखाना
अनिवार्य है कि कुछ फलनों की वह संगणना नहीं कर सकता या कुछ समस्याओं का
वह निर्णय नहीं कर सकता। उदाहरणतः हम दिखा सकते हैं कि प्रत्येक फलन की
संगणना ट्यूरिंग मशीनों द्वारा नहीं की जा सकती और प्रत्येक समस्या का
निर्णय ट्यूरिंग मशीनों द्वारा नहीं किया जा सकता। तब चर्च--ट्यूरिंग
अभिधारणा का आह्वान करके निष्कर्ष निकाल सकते हैं कि न केवल ट्यूरिंग
मशीनें प्रत्येक फलन की संगणना करने के लिए पर्याप्त शक्तिशाली नहीं हैं,
बल्कि कोई भी प्रभावी प्रक्रिया ऐसा नहीं कर सकती।

ऐसे निषेधात्मक परिणाम सिद्ध करने की कुंजी यह तथ्य है कि हम स्वयं
ट्यूरिंग मशीनों को संख्याएँ दे सकते हैं। ऐसा करने का सबसे सरल उपाय उनका
परिगणन करना है: संभवतः ट्यूरिंग मशीनों और उनके प्रोग्रामों को लिखने की
कोई विशिष्ट रीति निश्चित करके, फिर उन्हें क्रमबद्ध ढंग से सूचीबद्ध
करना। एक बार यह देख लेने पर कि ऐसा किया जा सकता है, ट्यूरिंग-असंगणनीय
फलनों का अस्तित्व समुच्चयों के आकार संबंधी सरल विचारों से निकलता है:
$\Nat$ से~$\Nat$ तक के फलनों का समुच्चय (वस्तुतः केवल $\Nat$ से
$\{0, 1\}$ तक के फलनों का भी) !!{nonenumerable} है, किंतु चूँकि हम सभी
ट्यूरिंग मशीनों का परिगणन कर सकते हैं, इसलिए ट्यूरिंग-संगणनीय फलनों का
समुच्चय केवल !!{denumerable} है।

हम ऐसे \emph{विशिष्ट} फलन और समस्याएँ भी परिभाषित कर सकते हैं जिन्हें
क्रमशः असंगणनीय और अनिर्णेय सिद्ध किया जा सके। ऐसी एक समस्या तथाकथित
\emph{विराम समस्या} है। ट्यूरिंग मशीनों के निर्देशों को सूचीबद्ध करके
उनका परिमित वर्णन किया जा सकता है। ट्यूरिंग मशीन का ऐसा वर्णन, अर्थात्
ट्यूरिंग मशीन प्रोग्राम, निश्चय ही किसी दूसरी ट्यूरिंग मशीन के आगत के
रूप में प्रयुक्त किया जा सकता है। अतः हम उन ट्यूरिंग मशीनों पर विचार कर
सकते हैं जो दूसरी ट्यूरिंग मशीनों के विषय में प्रश्नों का निर्णय करती
हैं। एक विशेष रूप से रोचक प्रश्न है: ``क्या दी गई ट्यूरिंग मशीन आगत~$n$
पर आरंभ किए जाने पर अंततः रुकती है?'' यदि इस प्रश्न का निर्णय करने वाली
ट्यूरिंग मशीन होती तो अच्छा होता: उसे ऐसी गुणवत्ता-नियंत्रण ट्यूरिंग
मशीन समझें जो सुनिश्चित करती हो कि ट्यूरिंग मशीनें अनंत चक्रों आदि में
न फँसें। ट्यूरिंग द्वारा सिद्ध रोचक तथ्य यह है कि ऐसी ट्यूरिंग मशीन हो
ही नहीं सकती। कोई एक ट्यूरिंग मशीन ऐसी नहीं हो सकती जो ट्यूरिंग मशीन
$M$ के वर्णन और किसी संख्या~$n$ से बने आगत पर आरंभ होकर सदैव निर्गत $1$
या $0$ के साथ इस अनुसार रुके कि $M$ आगत $n$ पर आरंभ होने पर रुकती या
नहीं।

विशिष्ट अनिर्णेय समस्याओं के उदाहरण मिल जाने पर हम उनका उपयोग यह दिखाने
के लिए कर सकते हैं कि अन्य समस्याएँ भी अनिर्णेय हैं। उदाहरणतः एक प्रसिद्ध
अनिर्णेय समस्या यह प्रश्न है, ``क्या प्रथम-कोटि !!{formula}~$!A$ वैध
है?'' ऐसी कोई ट्यूरिंग मशीन नहीं है जिसके बारे में यह सुनिश्चित हो कि
प्रथम-कोटि !!{formula}~$!A$ को आगत के रूप में दिए जाने पर वह इस अनुसार
निर्गत $1$ या $0$ के साथ रुकेगी कि $!A$ वैध है या नहीं। ऐतिहासिक रूप
से, इस समस्या को प्रभावी ढंग से हल करने वाली प्रक्रिया खोजने का प्रश्न
केवल ``निर्णय समस्या'' कहलाता था; अतः हम कहते हैं कि निर्णय समस्या
असमाधेय है। ट्यूरिंग और चर्च ने लगभग एक ही समय पर स्वतंत्र रूप से इस
परिणाम को सिद्ध किया था, इसलिए इसे चर्च--ट्यूरिंग प्रमेय भी कहते हैं।

\end{document}
